\section{Dropout 正则化}

\subsection{Dropout 的基本原理}

\begin{figure}[H]
\centering
\includegraphics[width=0.95\textwidth]{slides-images/slide-010.jpg}
\caption{Dropout：在训练时随机将部分激活值置零\protect\footnotemark}
\end{figure}
\footnotetext{Slides 第11页。}

Dropout 是一种正则化技术，核心思想是：\textbf{在训练时随机``关闭''一部分神经元，在测试时使用全部神经元}。

具体操作：
\begin{itemize}
    \item 训练时：以概率 $p$（通常 $p = 0.5$）将每个激活值随机设为 0
    \item 测试时：使用所有激活值，但乘以 $p$ 以保持期望值一致
\end{itemize}

\begin{lstlisting}[caption={Dropout 的 Python 实现示意},language=Python]
# Training
mask = (np.random.rand(*H.shape) > p)  # dropout mask
H_train = H * mask                     # zero out

# Testing
H_test = H * p                         # scale by p
\end{lstlisting}

\subsection{Dropout 为什么有效？}

\begin{figure}[H]
\centering
\includegraphics[width=0.85\textwidth]{slides-images/slide-011.jpg}
\caption{Dropout 的直觉解释：迫使模型学习冗余表示，不能过度依赖特定特征\protect\footnotemark}
\end{figure}
\footnotetext{Slides 第12页。}

\begin{importantbox}{Dropout 有效的直觉}
考虑一个猫分类器：模型可能学到``有耳朵 + 有毛皮 = 猫''这样的模式。但如果训练时``耳朵''特征有 50\% 概率被丢弃，模型就\textbf{不能过度依赖}耳朵特征，必须同时学会利用尾巴、爪子、毛皮等其他特征。

这迫使模型学习\textbf{冗余}和\textbf{分布式}表示：
\begin{itemize}
    \item 训练时表现略差（因为只看到部分特征）
    \item 测试时表现更好（因为学到了更鲁棒的表示）
\end{itemize}
\end{importantbox}

\begin{warningbox}{Dropout 的测试时缩放}
如果训练时 $p = 0.5$（50\% 概率保留），测试时使用全部激活值，那么每个神经元接收的输入总和会\textbf{翻倍}。为了保持训练和测试行为一致，测试时需要将所有激活值乘以 $p$。

不做这个缩放会导致测试时输出值远大于训练时，产生严重的行为偏差。
\end{warningbox}

\subsection{Dropout 的反向传播}

Dropout 的反向传播非常简单：被置零的激活值梯度也为零（类似 ReLU 的负半轴），不参与权重更新。这意味着每次训练迭代中，只有一部分参数被更新。

\begin{knowledgebox}{Dropout 与集成学习的联系}
Dropout 可以看作一种\textbf{隐式集成}：每次前向传播都在训练一个不同的``子网络''（由存活的神经元构成）。测试时使用全部神经元，相当于对所有可能的子网络进行了（近似的）集成。集成学习是机器学习中经典的提升泛化能力的方法。
\end{knowledgebox}

\subsection{本章小结}

Dropout 通过在训练时引入随机性来防止过拟合。它的核心机制是：训练时随机丢弃部分神经元（增加难度），测试时使用全部神经元并缩放（获得更好泛化）。Dropout 主要用于全连接层；在现代 CNN 和 Transformer 中，其他正则化方法（如数据增强、权重衰减）更为常用。

%% ========== Part 4: 激活函数 ==========

